\documentclass[conference,letterpaper]{IEEEtran}
\IEEEoverridecommandlockouts
\usepackage[T1]{fontenc}
\usepackage[utf8]{inputenc}
\usepackage{newtxtext,newtxmath}
\usepackage{booktabs,tabularx,array,amsmath,multirow}
\usepackage{balance}
\usepackage{url}
\usepackage{xcolor}
\usepackage[hidelinks]{hyperref}
\usepackage{microtype}
\newcommand{\stE}{\textsf{E}}
\newcommand{\stP}{\textsf{P}}
\newcommand{\stU}{\textsf{U}}
\newcommand{\evid}{\mathcal{E}}
\renewcommand{\arraystretch}{1.11}
\begin{document}
\title{From AI Strategy to University Rules in the GCC}
% Anonymized research copy. Authors, acknowledgements, and camera-ready copyright deliberately excluded.
\author{}
\maketitle
\begin{abstract}
National artificial-intelligence (AI) strategies do not themselves establish public rules that university users can identify and apply. We audit a bilingual Arabic--English, four-layer documentary profile across all six Gulf Cooperation Council (GCC) states: national direction (L1), operational governance (L2), higher-education (HE) governance translation (L3), and institution-level rules (L4). The frozen, 23 July 2026 evidence frame comprises 36 official-source records. Under the primary substantive-control gate, L4 is explicit in all six jurisdictions while L3 is explicit in one; the gap remains in four jurisdictions under a broader interpretation of a recovered UAE HE evaluation framework. A corrected, pre-outcome oldest-four-with-ties corroboration panel covers 57 institutions (20 explicit, 12 partial, 25 unlocated). We distinguish documentary co-presence from demonstrated instrument-to-instrument transmission, document asymmetric L3/L4 qualification, and provide exact finite-corpus robustness bounds. A secondary interval-valued fuzzy representation of partial evidence retains a minimum L4--L3 envelope gap of 0.5 but cannot correct missing sources or coder error. The result is a reusable public-traceability audit, not a country ranking, prevalence estimate, compliance audit, or causal account of policy adoption.
\end{abstract}
\begin{IEEEkeywords}
AI governance, Gulf Cooperation Council, higher education, bilingual policy audit, documentary traceability, evidence admissibility, interval-valued uncertainty.
\end{IEEEkeywords}

\section{Introduction}
Generative AI (GenAI) raises immediate university governance questions: which uses are allowed in assessment, what must students disclose, which data may be submitted to external tools, and who remains accountable for verification? An aspirational national strategy does not answer these questions for an instructor or a student. The operational issue is whether the relevant instrument and its authority can be \emph{located in public evidence} at successive governance levels.

Comparative studies describe national strategies in the GCC~\cite{albous}, policies in 40 universities across six regions~\cite{jin}, 116 US research universities~\cite{mcdonald}, and 343 institutions across five countries~\cite{parker}. Assessment guidance~\cite{moorhouse}, comparative university frameworks~\cite{li}, and systematic/institutional reviews~\cite{jiang,khoza,an,bacalso,qian} reveal heterogeneous approaches. Their predominant unit of analysis is a policy or a university; our primary unit is a \emph{four-layer public-document profile within a jurisdiction}. This does not presume that a document in one layer caused an instrument in another. Governance is multi-level and multi-actor~\cite{vukasovic}, and policy-translation studies warn against confusing similar visible products with a documented transmission process~\cite{donina}.

We ask: (RQ1) Under an identical evidence gate and cutoff, which layers have explicitly identifiable instruments? (RQ2) Does the L3--L4 documentary discontinuity survive source recovery, a deterministic institutional corroboration panel, and transparent alternative representations of partial evidence? (RQ3) Which parts of an asserted vertical \emph{policy chain} are observable as document availability, and which would additionally require explicit citations or implementation data?

The contributions are threefold. First, a replicable categorical audit separates \emph{instrument existence}, \emph{documentary linkage}, and \emph{implementation}, which earlier policy-translation language can conflate. Second, a fixed-rule sample of 57 institutions tests whether the jurisdictional L4 existence finding rests on a single example. Third, exact, small-sample sensitivity bounds show which interpretation remains invariant while naming what those bounds \emph{do not} establish. We retain interval-valued fuzzy notation only as a concise representation of partial-state possibilities, not a ranking or discovery algorithm.

\begin{table}[!b]
\centering\caption{Comparison with adjacent research (object, not quality ranking)}\label{tab:related}
\footnotesize
\begin{tabularx}{\columnwidth}{@{}p{0.18\columnwidth}p{0.21\columnwidth}X@{}}
\toprule
Study & Coverage & Primary object\tabularnewline\midrule
Jin et al.~\cite{jin} & 40 HEIs, six regions & Institution policies and adoption\tabularnewline
McDonald et al.~\cite{mcdonald} & 116 US R1 HEIs & Institution guidance\tabularnewline
Parker et al.~\cite{parker} & 343 HEIs, five countries & Cross-country university policies\tabularnewline
Khoza et al.~\cite{khoza} & 35 HEIs, 43 documents & Institutional governance structures\tabularnewline
This audit & Six GCC states; 57 HEIs & Vertical L1--L4 public evidence; deterministic corroboration\tabularnewline
\bottomrule
\end{tabularx}
\end{table}

\section{Construct and Source-Admissibility Protocol}
\subsection{Different observable claims}
For jurisdiction $j$ and layer $\ell$, define the categorical documentary state
\begin{equation}\label{eq:chain}
 c_j=(L_{1j},L_{2j},L_{3j},L_{4j}),\quad L_{\ell j}\in\{\stE,\stP,\stU\}.
\end{equation}
An explicit state $\stE$ means an identifiable, inspectable official instrument meets the \emph{layer-specific} gate. A partial state $\stP$ records an official but narrower, indirect, developmental, or otherwise non-qualifying lead. An unlocated state $\stU$ means no qualifying instrument was located under the documented bilingual public-search closure; $\stU$ neither asserts that no internal rule exists nor alleges noncompliance. These are nominal labels, not ordinal country ratings.

The following three claims must not be interchanged. \emph{Availability} asks whether a qualifying public instrument exists at L1, L2, L3, or L4. \emph{Linkage} requires an inspectable instrument-to-instrument reference, delegation, scope map, or expressly assigned responsibility connecting particular levels. \emph{Implementation} requires evidence of actual institutional use, enforcement, monitoring, or participant experience. The corpus codes availability only. Even a fully $\stE$--$\stE$--$\stE$--$\stE$ vector is therefore \emph{not} proof of causal transmission or administrative alignment. A future link ledger would code directed, dated instrument references as separate edges, never infer them from co-presence.

\begin{table}[t]
\centering\caption{Layer-specific explicit gates: different scope, same evidential discipline}\label{tab:gates}
\footnotesize
\begin{tabularx}{\columnwidth}{@{}p{0.11\columnwidth}X@{}}
\toprule
Layer & Requirement for $\stE$ (not merely an announcement)\tabularnewline\midrule
L1 & Adopted or authoritative national AI direction with identifiable issuer.\tabularnewline
L2 & Inspectable, cross-sector operational guideline, control, regulation, or equivalent.\tabularnewline
L3 & Authoritative HE-specific, or education-sector instrument explicitly applicable to HE, \emph{with substantive AI-governance controls}.\tabularnewline
L4 & At least one applicable university/college AI rule with institutional authority and substantive scope.\tabularnewline
\bottomrule
\end{tabularx}
\end{table}

\subsection{The intentionally asymmetric L3 and L4 estimands}
An L3-positive classification is a \emph{sector-level documentary proposition}. For $\stE$, an HE-applicable national or education-sector instrument must specify meaningful controls (e.g., assessment, disclosure, privacy, accountability, or oversight). School-only digital modernisation, a ministry working group still preparing guidance, and a framework measuring AI readiness without prescribing use controls do not satisfy this gate. The resulting $\stP$ or $\stU$ classification follows its documented scope and evidence status.

An L4-positive classification is a much narrower \emph{existential} proposition: \emph{one} qualifying institution in jurisdiction $j$ has an applicable, authoritative public rule. It is not a claim that all, most, or a representative fraction of the country's universities have one. Accordingly, the contrast $L3_j\ne\stE$, $L4_j=\stE$ is a discontinuity between two specific public-evidence propositions. It is neither a ratio of comparable populations nor a gap in verified implementation rates. The panel below provides separate finite-case context, not an estimator of sector-wide prevalence.

\subsection{Frozen source frame and bilingual procedure}
The analytical cutoff is 23 July 2026. The 36-record official register includes 7 L1, 8 L2, 5 L3, 11 L4, and 5 GCC-regional records. Nineteen were located in Arabic and English, 15 in English only, and two in Arabic only; alternative language versions of one official item do not count as separate instruments. Official government, regulator, ministry and institutional sources take precedence over commentary and search snippets. A dated official announcement can establish \emph{development} or \emph{approval}, but cannot substitute for the text of a required operative instrument.

The search protocol includes Arabic and English families for AI/GenAI, strategy, regulation, policy, governance, education/HE, teaching, assessment, integrity, responsibility, disclosure, privacy, and oversight, combining terms with official domains and document libraries. For each candidate, the audit fields are source ID, jurisdiction, layer, issuing authority, exact title and edition, issue/effective and retrieval dates, language, direct official locator, provision/page, proposed state, and explicit reason for meeting or failing that layer's gate. A $\stU$ closure additionally records searches of the relevant institutional policy library, handbook, integrity/teaching pages and official-domain Arabic/English results. Changing web architecture and search visibility can change $\stU$ without a real policy change; subsequent recovery of an eligible \emph{pre-cutoff} document is logged, not retroactively presented as already discovered.

Arabic/English parity requires checking issuer, scope, approved version, effective date, authority hierarchy, and controlling-language statements separately from translation quality. Qatar University's instructions, for example, identify the English text as a translation for which Arabic prevails when inconsistent~\cite{qataru}. If the authoritative text or a decisive discrepancy is inaccessible, the classification must remain provisional. The current study is not a completed independent bilingual adjudication experiment: no second-coder agreement statistic is claimed.

\subsection{Recovery and the unchanged admission gate}
Pre-cutoff source recovery found Sultan Qaboos University's 2024 teaching/learning/assessment GenAI guidance~\cite{squ}, making Oman's L4 $\stE$; omitting this recovered document yields a separately reported sensitivity. The UAE's 23 March 2026 HE Outcome-Based Evaluation Framework (OBEF) v11.5~\cite{obef} contains AI-enabled educational design, assessment/analytics, and institutional-readiness dimensions. Under the \emph{primary control gate}, this evaluation framework is $L3=\stP$: a criterion for evaluating AI integration is not itself an operative responsible-use control. A deliberately broader sensitivity treats it as $\stE$. The comparison tests gate dependence without silently redefining L3 after seeing the result.

\section{An Operational Traceability Contract}
\subsection{Why an apparent vertical chain needs two additional tests}
The intuitive phrase ``from strategy to university rules'' may be read in three ways. The weakest, investigated here, is that suitable public instruments can be found at several governance levels. A stronger documentary statement is that an identifiable lower-level instrument \emph{expressly draws authority or control requirements} from a named upper-level one. Stronger still, actual teaching, assessment or data-handling behaviour might conform to those rules. These are distinct measurement targets, even when their narrative descriptions are similar.

For each jurisdiction and layer, let $A_{\ell j}$ denote the proposition that an authority- and scope-qualified public instrument is located. Then $L_{\ell j}=\stE$ is a recorded positive witness for $A_{\ell j}$ under Table~\ref{tab:gates}; $\stP$ and $\stU$ encode incomplete \emph{evidence}, not the falsity of the proposition. For an adjacent pair of layers define a separate documentary-link proposition $R_{\ell,\ell+1,j}$, requiring an identifiable source-to-target relationship, e.g., a lower instrument that explicitly invokes the upper instrument, a dated delegation, or a controlling sector circular with named institutional applicability. Define an implementation proposition $I_{\ell j}$ requiring records of use, compliance or enforcement in the relevant institutional setting. Our reported chain outcomes concern $A$; the $R$ and $I$ variables were not systematically examined or coded. They are \emph{unassessed}, not coded ``no'' and not automatically inferred from $A$.

This separation matters even for Saudi Arabia's fully explicit four-layer vector. Four admissible instruments do not demonstrate three directed links. Conversely, an L3 source unlocated in this public search does not rule out a substantive implementation pathway through unpublished ministerial communications. A later study could extend, rather than overwrite, the present frame by building an evidence graph with instrument-version nodes and directed reference edges. Each claimed edge should carry both the source and target ID, a quoted authority/obligation passage, the exact language/edition and the date when the edge was operative. Without this information, a schematic arrow is illustrative rather than an observed transmission fact.

\begin{table*}[t]
\centering\caption{Evidence contract separating available instruments, explicit documentary links, and implementation. Only the first object is measured in the present GCC results.}\label{tab:contract}
\footnotesize
\begin{tabularx}{\textwidth}{@{}p{0.16\textwidth}p{0.29\textwidth}p{0.29\textwidth}X@{}}
\toprule
Object & Minimum positive witness & What the witness supports & What it cannot independently establish\tabularnewline\midrule
Layer instrument $A_{\ell j}$ & Official, dated, inspectable instrument meeting the layer's authority and scope gate & \emph{Public existence} at that layer, coded $\stE/\stP/\stU$ & Explicit transmission to an adjacent document or any downstream behavior\tabularnewline
Directed link $R_{\ell,\ell+1,j}$ & Instrument-to-instrument citation, directive, delegation or explicit application clause & A specific \emph{documentary reference edge} between named instrument versions & Effective compliance or its substantive impact; \emph{not coded here}\tabularnewline
Implementation $I_{\ell j}$ & Auditable enforcement records, observed assessment practice, records review, or field interviews & Some \emph{realized rule operation} in a defined setting/time period & Nation-wide causal adoption; \emph{not observed here}\tabularnewline
\bottomrule
\end{tabularx}
\end{table*}

\subsection{Admission is not equivalent across layers}
A comparison is admissible only after naming its estimand. At L3, the question is whether a sector-applicable governance instrument exists with content that actually controls AI use in higher education. At L4, the question is whether \emph{any one} institution has issued such a rule for its own users. The former is a jurisdictional \emph{HE-sector instrument}; the latter a jurisdictional \emph{existence witness}. Replacing $A_{3j}$ with the count of local university documents would make the original gap disappear by definition while answering a different research question. Likewise, insisting that all universities must publish a rule to earn L4 $\stE$ would create a prevalence gate for which the study never assembled a representative denominator. The asymmetric gate is thus purposeful but must be stated at every interpretation of the L3--L4 contrast.

The approach is designed to be conservative about instrument authority while permissive about discovery. A direct official policy in one language is retained even if there is no translation; a second-language mirror cannot count as a second case. Where an official news article reports a proposed policy, it can be admitted as process evidence $\stP$ but does not receive the status of a promulgated control document. A course syllabus with a ChatGPT prohibition may evidence a binding \emph{course} rule, yet is not automatically an institution-wide GenAI governance instrument. The policy question becomes more precise, not merely more demanding: the coder records exactly \emph{which scope} the located document governs.

\section{Pre-Outcome Institutional Corroboration}
\subsection{Sampling identity and correction trail}
The L4 jurisdictional gate can close with one institution. For corroboration, a distinct panel was specified \emph{before} examining AI-policy outcomes. Eligible officially recognized or licensed civilian HEIs in each state were separated into G1 public/general, G2 private/non-government and G3 public applied/technical/specialized. Within each jurisdiction--stratum cell, the oldest four according to the current institution's official founding year were selected, ties at fourth place retained, and smaller cells fully included. Institutions directly affiliated with authors were excluded before policy search. Founding year is merely an auditable ordering field: it implies neither seniority-related quality nor policy maturity.

Two retrospective corrections were necessary for \emph{implementing} the unchanged rule. The originally omitted Kingdom University (Bahrain G2, founding year 2001) entered after a pre-cutoff regulator report and index supported its eligibility, removing two later-founded cases~\cite{bqa1,bqa2}. The UAE Commission for Academic Accreditation listed Emirates Institute for Banking and Financial Studies (EIBFS) as an active HEI founded in 1983; it entered UAE G2 and displaced two 1997 cases below the now-corrected threshold~\cite{caa1,caa2}. EIBFS was \emph{then} searched under the same AI-policy gate and was $\stU$. Membership selection and outcome coding were not retrofitted to improve results. The available retrospective record does not, however, reconstruct every excluded eligible HEI; hence the analysis makes no claim of a complete oldest-five candidate universe.

\subsection{Institutional classification and boundary}
The corrected panel contains $N=57$ institutions: UAE 13, Saudi Arabia 11, Qatar 10, Oman 8, Bahrain 7 and Kuwait 8, with G1/G2/G3 totals 14/24/19. An institutional $\stE$ requires an inspectable AI/GenAI rule with substantively applicable institutional authority; course-only use restrictions, research announcements, platform-limited clauses and draft proposals are generally $\stP$ when they contain direct relevant language, not $\stE$ by association. With no qualifying pre-cutoff material after documented bilingual domain-and-library closure, the state is $\stU$; silence is never coded as explicit absence. The companion 57-institution key retains all clickable official evidence/search locators and distinguishes a context/search homepage from direct rule evidence.

\begin{table*}[t]
\centering\caption{Nominal public-evidence profiles at the fixed cutoff (black-and-white legible; $\stE$=explicit, $\stP$=partial, $\stU$=unlocated). \emph{No adjacency or causal transmission is inferred.}}\label{tab:profiles}
\footnotesize
\begin{tabularx}{\textwidth}{@{}p{0.16\textwidth}ccccX@{}}
\toprule
Jurisdiction & L1 & L2 & L3 & L4 & Principal interpretation\tabularnewline\midrule
UAE & \stE & \stE & \stP & \stE & HE OBEF is assessment/readiness evaluation, not an operative HE control under the primary gate.\tabularnewline
Saudi Arabia & \stE & \stE & \stE & \stE & All four public-instrument gates satisfied; still no demonstrated directed transfer.\tabularnewline
Qatar & \stE & \stE & \stP & \stE & HE-specific public governance bridge is below the substantive-control gate.\tabularnewline
Oman & \stE & \stE & \stU & \stE & Recovered 2024 SQU teaching/assessment guidance qualifies as L4.\tabularnewline
Bahrain & \stP & \stE & \stU & \stE & School-directed digital-empowerment program does not qualify as HE-specific L3.\tabularnewline
Kuwait & \stP & \stP & \stU & \stE & A formal Kuwait University GenAI rule is identifiable despite partial upstream instruments.\tabularnewline
\bottomrule
\end{tabularx}
\end{table*}

\section{Results and Exact Sensitivities}
\subsection{Layer-by-layer outcome}
Table~\ref{tab:profiles} makes the central contrast inspectable without relying on a color heatmap. The counts are L1 $4\stE+2\stP$, L2 $5\stE+1\stP$, L3 $1\stE+2\stP+3\stU$, and L4 $6\stE$. Define
\begin{equation}\label{eq:discontinuity}
D=\sum_{j=1}^{6}\mathbf{1}\{L_{4j}=\stE\ \wedge\ L_{3j}\neq\stE\}=5.
\end{equation}
Thus five jurisdictions show explicit institution-level public rules with no primary-gate explicit sector-specific HE control document; Saudi Arabia is the sole full $\stE$--$\stE$--$\stE$--$\stE$ profile. $D$ describes \emph{co-presence across differently scoped documentary gates}, not a count of broken regulatory implementation pathways.

\begin{table}[t]
\centering\caption{Finite-corpus source and gate sensitivity (not confidence intervals)}\label{tab:robustness}
\footnotesize
\begin{tabularx}{\columnwidth}{@{}Xccc@{}}
\toprule
Condition & L3 $\stE$ & L4 $\stE$ & $D$\tabularnewline\midrule
Pre-recovery SQU omission & $1/6$ & $5/6$ & $4/6$\tabularnewline
Primary recovered evidence & $1/6$ & $6/6$ & $5/6$\tabularnewline
UAE OBEF broader gate & $2/6$ & $6/6$ & $4/6$\tabularnewline
Leave Saudi Arabia out & $0/5$ & $5/5$ & $5/5$\tabularnewline
Leave any other state out & $1/5$ & $5/5$ & $4/5$\tabularnewline
\bottomrule
\end{tabularx}
\end{table}

A monotone ``staircase'' story would predict $L4_j=\stE$ only where every preceding layer is explicit. Five vectors contradict that proposition, and four still contradict it under broader UAE admission. This rejects \emph{public-documentary monotonicity} within the six observed vectors, but does not establish the temporal origin of university rules. Ministries may issue non-public guidance, institutions may exercise legitimate autonomy, and public web archives may omit existing instruments.

\subsection{Partial-state sensitivity, without artificial precision}
Let $n_{\ell,\stE}$ and $n_{\ell,\stP}$ count explicit and partial states in layer $\ell$. The interval-valued fuzzy \emph{representation} for the sole proposition ``explicitly traceable public instrument'' maps $\stE\mapsto[1,1]$, $\stP\mapsto[0,1]$, and $\stU\mapsto[0,0]$, yielding
\begin{equation}\label{eq:fuzzy}
\widetilde F_\ell=\left[\frac{n_{\ell,\stE}}6,\ \frac{n_{\ell,\stE}+n_{\ell,\stP}}6\right].
\end{equation}
These are elementary lower/upper envelopes implied by the categorical counts, consistent with interval-valued fuzzy-set notation~\cite{petry,pekala}; they are not model-estimated membership functions. Substitution gives $\widetilde F_1=[.667,1]$, $\widetilde F_2=[.833,1]$, $\widetilde F_3=[.167,.500]$, and $\widetilde F_4=[1,1]$. Because even both L3 partial records cannot yield more than $3/6$ explicit membership, every allowed realization satisfies
\begin{equation}\label{eq:bound}
F_4-F_3\geq 1-\frac{1+2}{6}=\frac12.
\end{equation}
The deliberate OBEF reclassification moves an L3 $\stP$ to $\stE$ but leaves its upper envelope $3/6$ and the minimum gap unchanged. This bound merely makes the worst permitted treatment of $\stP$ transparent; it neither identifies omitted sources nor quantifies coder uncertainty, search recall, quality, or prevalence. The categorical states, not fuzzy averages, remain the primary empirical results.

An additional deterministic \emph{conditional} stress test avoids pretending that agreement was measured. If at most $k$ jurisdictional indicator outcomes in Eq.~\eqref{eq:discontinuity} are changed by an adversarial coding or source-admission revision, then $D_{\rm revised}\geq\max(0,5-k)$. In particular, the direction $D_{\rm revised}>0$ persists when $k\leq4$. No evidence is supplied that $k$ is small, and this algebra cannot substitute for independent recoding. The contribution is a clearly stated robustness threshold, not a probability of correctness.

\subsection{The fixed institutional panel, not a population rate}
\begin{table*}[t]
\centering\caption{Corrected 57-institution panel: $(\stE,\stP,\stU)$ in each pre-outcome jurisdiction--stratum cell. Unequal cell sizes are intentional.}\label{tab:institution}
\footnotesize
\begin{tabularx}{\textwidth}{@{}lcccc@{}}
\toprule
Jurisdiction & G1 public/general & G2 private/non-government & G3 public applied/specialized & Total\tabularnewline\midrule
UAE & $3,0,1$ (4) & $1,2,1$ (4) & $1,1,3$ (5) & $5,3,5$ (13)\tabularnewline
Saudi Arabia & $2,1,1$ (4) & $1,1,2$ (4) & $0,2,1$ (3) & $3,4,4$ (11)\tabularnewline
Qatar & $1,0,1$ (2) & $1,2,1$ (4) & $0,0,4$ (4) & $2,2,6$ (10)\tabularnewline
Oman & $1,0,0$ (1) & $1,1,2$ (4) & $0,0,3$ (3) & $2,1,5$ (8)\tabularnewline
Bahrain & $1,0,0$ (1) & $3,0,1$ (4) & $2,0,0$ (2) & $6,0,1$ (7)\tabularnewline
Kuwait & $1,0,1$ (2) & $1,2,1$ (4) & $0,0,2$ (2) & $2,2,4$ (8)\tabularnewline\midrule
All six & $9,1,4$ (14) & $8,8,8$ (24) & $3,3,13$ (19) & $20,12,25$ (57)\tabularnewline
\bottomrule
\end{tabularx}
\end{table*}

Table~\ref{tab:institution} reports all fixed-cell outcomes. Each jurisdiction contains at least two $\stE$ and at least one non-$\stE$ institutional record. The $20$ explicitly coded rules corroborate L4 existence without supporting universal coverage; $12$ partially supported and $25$ unlocated cases instead expose heterogeneity in this non-random bounded selection. Because ties and small strata lead to unequal cells, raw shares are $20/57=.351$, $12/57=.211$, and $25/57=.439$, whereas equal jurisdiction--stratum-cell weighting changes them to $.414$, $.173$, and $.413$. An oldest-three-with-ties sensitivity retains $47$ institutions ($19\stE$, $6\stP$, $22\stU$). These changes are sensitivity to the \emph{sampling design}, not estimates of population percentages, bias corrections, or country league tables.

\section{Interpretive Stress Tests}
\subsection{Alternative explanations are not the same as contradictions}
The observed $D=5$ is compatible with several different processes. A university can act through academic-integrity authority before a ministry publishes a cross-sector HE instrument; a ministry may communicate requirements without an accessible circular; or a binding document may be stored on an unindexed policy portal. Any of these can coexist with the locked $\stE/\stP/\stU$ evidence matrix. Distinguishing them requires data outside an official-web documentary search. Our result therefore rejects a literal claim that \emph{all four layers must be visibly explicit whenever any university rule is visible}; it does not reject coordinated government guidance, local compliance, or policy diffusion as real-world mechanisms.

Three counterfactuals clarify the observed versus non-observed margins. First, suppose a new valid, pre-cutoff HE governance control for Oman is recovered. The primary matrix and the published coding decision would need a transparent source-admission revision: Oman would change $L3:\stU\rightarrow\stE$ and $D$ would fall by one. Second, suppose an already indexed policy merely includes a new summary sentence mentioning AI but provides no control or authority. It would be $\stP$, not a valid promotion to $\stE$. Third, suppose an identified university policy explicitly cites its national strategy. That new evidence would support a specific directed-link observation $R_{1,4,j}$ but leave the existing four-layer availability vector unchanged. These experiments illustrate how the construct prevents evidence from being counted at the wrong level.

\subsection{A typed error register rather than one uncertainty score}
The sources of uncertainty are heterogeneous and cannot be collapsed into a single fuzzy membership interval. \emph{Source recovery} concerns missed qualifying instruments; \emph{source identity} concerns edition and effective date; \emph{semantic coding} concerns whether a located clause meets a control gate; \emph{language authority} concerns which version governs when bilingual texts diverge; \emph{selection-boundary error} changes panel membership; and \emph{web disappearance} affects reproducibility of a locator. These mechanisms demand different controls: bilingual repeat search, signed source snapshot, blind recoding, recorded language precedence, official eligibility ledger, and immutable archiving, respectively. The $\stP\mapsto[0,1]$ interval only represents one declared \emph{interpretation range} for already located partial evidence. It is not a sensitivity analysis over all possible missing evidence or all human judgements.

\begin{table}[t]
\centering\caption{Threat-to-evidence mapping for audit replication}\label{tab:threats}\footnotesize
\begin{tabularx}{\columnwidth}{@{}p{0.33\columnwidth}X@{}}
\toprule
Threat & Required independent control\tabularnewline\midrule
Unseen official document & Replicate Arabic/English search, register recovery\tabularnewline
Wrong edition or date & Record document version, cutoff, source hash\tabularnewline
Ambiguous control scope & Blind second coding of clause and gate\tabularnewline
Translation/authority conflict & Read controlling language; log adjudication\tabularnewline
Wrong sample boundary & Freeze regulator eligibility and founding-year list\tabularnewline
Broken/moved URL & Archive source snapshot and capture time\tabularnewline
\bottomrule
\end{tabularx}
\end{table}

\subsection{What would overturn the narrower finding?}
For the primary categorical result to be eliminated entirely, every one of the five jurisdictions currently meeting $L4=\stE,L3\ne\stE$ would have to lose its L4 explicit witness or acquire a qualifying L3 explicit HE instrument (or a combination of such changes). This is the content of $D\geq5-k$, not evidence that five independent coding errors are improbable. If only one decisive record is disputed, the matrix difference is at most one jurisdictional indicator in this specific primary result. This claim remains conditional on the six official-source profiles and does not prevent a later, better documented corpus from substantially revising the map.

The broader UAE-OBEF interpretation is informative because it deliberately expands what may count as sector translation but still leaves three L3 unlocated jurisdictions. It is \emph{not} a test of implementation, and the fact that the partial-state interval upper bound stays fixed does not show that the policy content is invariant. The two kinds of sensitivity answer separate questions: alternative L3 gate definitions change a nominal judgement; the interval envelope allows any membership within partial cases given the observed judgement. Reporting both and calling neither a sampling confidence interval makes the findings more reusable across new evidence windows.

\section{Decisive Record Examples and Transfer Rules}
\subsection{What a verifiable document-level witness looks like}
Table~\ref{tab:anchors} gives a small, non-exhaustive set of documentary anchors behind decisive boundary choices. The links are to official authorities or institutions retained in the original evidence lineage; they should be opened and version-checked at deposition. They are not a substitute for release of the complete 36-row register. The six institutional examples include one positive L4 witness for every jurisdiction. A direct policy PDF is stronger for audit replay than a generic university homepage: an official landing page can locate a record but need not disclose a stable, signed policy edition. Conversely, a non-qualifying source is scientifically useful when its \emph{scope limitation} is precisely stated rather than discarded without explanation.

\begin{table*}[t]
\centering\caption{Selected official L4 anchors for cross-jurisdiction existence checking. Each locator is clickable; source editions and authority remain subject to the full ledger audit.}\label{tab:anchors}\footnotesize
\begin{tabularx}{\textwidth}{@{}p{0.13\textwidth}p{0.34\textwidth}p{0.30\textwidth}X@{}}
\toprule
Jurisdiction & Institution / documented object & Why it is an L4 witness under the locked gate & Official evidence locator\tabularnewline\midrule
UAE & United Arab Emirates University / student GenAI usage policy & Institution-authorized student and assessment use requirements & \href{https://www.uaeu.ac.ae/en/about/pdfs/policies-2025/generative-ai-policy.pdf}{UAEU policy PDF}\tabularnewline
Saudi Arabia & King Saud University / education GenAI guide & Official university guidance on AI in teaching and learning & \href{https://aio.ksu.edu.sa/ar/node/2976}{KSU AI Office page}\tabularnewline
Qatar & Qatar University / approved GenAI instructions & Institution-wide rules, disclosure and responsibility; authoritative Arabic version controls & \href{https://www.qu.edu.qa/en-us/about/policy/Documents/Instructions\%20onThe\%20Use\%20of\%20Generative\%20Artificial\%20Intelligence\%20Tools\%20at\%20Qatar\%20University-Approved.pdf}{Qatar University PDF}\tabularnewline
Oman & Sultan Qaboos University / 2024 GenAI teaching and assessment guidance & Recovered pre-cutoff institutional guidance; refer to specific archived edition & \href{https://www.squ.edu.om/About/Support-Centers-/Centre-For-Excellence-In-Teaching-and-Learning}{SQU teaching centre}\tabularnewline
Bahrain & University of Bahrain / AI usage regulation & University governance regulation with assessment, verification and accountability clauses & \href{https://governance.uob.edu.bh/view/297/entry/156/?pagenum=3}{Official governance registry}\tabularnewline
Kuwait & Kuwait University / formal GenAI policy & Public university-wide GenAI policy and related university approval & \href{https://www.ku.edu.kw/sites/default/files/2024-12/AI\%20Policy\%202024\%20November\%20-\%20English.pdf}{Kuwait University PDF}\tabularnewline
\bottomrule
\end{tabularx}
\end{table*}

A key counterexample to document-level overreach is the Bahrain Ministry of Education school-oriented Digital Empowerment Programme~\cite{bahrain}. Its visibility is not evidence of a higher-education instrument; promoting it to L3 $\stE$ would change the estimand. The UAE OBEF~\cite{obef}, by contrast, is truly higher-education-specific but emphasizes performance evaluation and AI integration rather than prescriptive AI-use controls. It therefore contributes directly relevant L3 $\stP$ evidence and motivates the broader admission sensitivity. This difference between \emph{scope mismatch} and \emph{control-function mismatch} is reusable in other multilingual governance audits: both may yield a non-explicit state, yet the next documentation search is different.

\subsection{How a subsequent researcher can improve rather than overwrite the record}
A longitudinal replication should assign a new evidence-window identifier, preserve the old document and classification digest, and log each newly found item as ``pre-cutoff recovery,'' ``post-cutoff adoption,'' ``edition replacement,'' or ``coding correction.'' These categories have different scientific implications. A recovered 2024 policy alters what the earlier public search located, not what the university had actually done in the meantime. A new 2027 policy is instead a change in the later documentary landscape and should not rewrite the 2026 matrix. An updated bilingual translation can clarify the controlling text without proving a new policy adoption event. A coders' disagreement changes the uncertainty about classification but is not itself an observed government action.

The same discipline applies to negative-search records. For $\stU$, a reviewer should be able to see the exact sites searched, query families in both languages, most recent page scan date, and a short reason that candidate pages failed the gate. Recording ``not found'' without those surfaces makes a zero impossible to distinguish from a missing search. To make the protocol transferable, future researchers should retain both the official short title and a normalized unique document identity across editions, together with a page/paragraph pointer to the control provisions. The evidence register is then usable as a versioned, auditable corpus rather than only a table of policy names.

\section{Auditable Source Boundaries and Practical Reuse}
\subsection{Evidence examples that control decisive inferences}
Several high-impact classifications illustrate why URL provenance is not interchangeable with an inferred chain edge. Sultan Qaboos University\textquotesingle s 2024 teaching/assessment guidance~\cite{squ} supports Oman L4, but not Oman L3. The UAE OBEF v11.5~\cite{obef} is a ministry-issued HE evaluation instrument, but a primary control threshold keeps its L3 state partial. Bahrain\textquotesingle s Digital Empowerment in Education Program~\cite{bahrain} is school-facing and cannot be upgraded to an HE policy. Qatar University\textquotesingle s approved instructions~\cite{qataru} provide applicable institutional instructions, but their bilingual hierarchy requires checking the authoritative Arabic version rather than assuming both languages are independently binding. No one of these cases proves that a lower-level rule was issued \emph{because of} a higher-level strategy.

\subsection{Eight fields for an inspectable public rule}
A practical outcome of the audit is the proposed documentation schema in Table~\ref{tab:eight}: its fields identify \emph{what must be public to check a claim} rather than offering a new quality ranking. The schema follows recurring documentary needs in the audited instruments and overlaps with UNESCO and OECD guidance~\cite{unesco,oecd}. That overlap is construct alignment, not independent validation of eight factors, an empirically optimized checklist, or proof that their presence makes an institution safe or compliant.

\begin{table}[t]
\centering\caption{Proposed documentation fields; none is a validated score}\label{tab:eight}
\footnotesize
\begin{tabularx}{\columnwidth}{@{}p{0.32\columnwidth}X@{}}
\toprule
Public field & Documentary question\tabularnewline\midrule
Authority and version & Who approved which operative text, and when?\tabularnewline
Scope and roles & Which people, programs, and activities are covered?\tabularnewline
Assessment use & What is permitted, conditional, or prohibited?\tabularnewline
Disclosure & What AI involvement requires acknowledgment?\tabularnewline
Data and tools & Which data classes and external tools are restricted?\tabularnewline
Human responsibility & Who verifies outputs and retains accountability?\tabularnewline
Fairness and challenge & Where are accessibility, bias and appeal addressed?\tabularnewline
Maintenance and review & Who monitors changes and publishes revisions?\tabularnewline
\bottomrule
\end{tabularx}
\end{table}

\subsection{A repeatable source-to-claim audit trail}
Independent replication needs more than hyperlinks to institutional homepages. Each 36-record source row must be delivered with a \emph{direct document locator}, recorded edition and retrieval date, quote or exact clause/page locator, primary coding gate, reason for partial/unlocated status, alternative-language checks, and stable capture/hash where redistribution is permitted. For the 57 institutions, the companion key provides the locked labels and clickable retained locators in a readable standalone format; a URL to an official homepage for an unlocated case remains a \emph{search locator}, not a supposedly qualifying policy document. A replay should first verify source identity and cutoff, then reconstruct the $\stE/\stP/\stU$ matrix, then regenerate Tables~\ref{tab:profiles}--\ref{tab:institution} and Eqs.~\eqref{eq:discontinuity}--\eqref{eq:bound}. Later web changes belong in a new dated audit window rather than an unreported alteration of the original results.

For prospective independent coding, two coders should independently receive exact frozen source snapshots, a gatebook and two blinded record orders. They should classify (i) source qualification by layer and (ii) any explicit cross-reference in separate fields, with source passages and Arabic/English authority decisions recorded before adjudication. Report a $3\!\times\!3$ nominal confusion matrix for $\stE,\stP,\stU$, raw agreement, and Cohen's $\kappa$ (plus case-specific disagreements), then run both unadjudicated and adjudicated chain results. This is a \emph{protocol for future validation}: independent double-coding, actual adjudication, and agreement estimates were \emph{not performed} for the present results. The source register and its licensed snapshots likewise require an explicit, versioned public deposition to meet full replication expectations; mentioning a ledger does not by itself prove public access.

\subsection{Minimal reproducibility objects and admission order}
The necessary deposit has three linked resources and a reconstruction script, not only a paper PDF. The first is a \emph{36-row official-source register}: source ID, layer, country, canonical document title, issuing authority, policy/status, precise issue/effective date and edition, controlling language, direct URL, quoted clause and page, source capture date, source file hash, E/P/U rationale, and post-cutoff exclusion flag. The second is an \emph{institutional eligibility ledger}: every candidate capable of altering the oldest-four boundary, regulator identity, founding year, stratum and exclusion or tie reason. The third is the \emph{57-row outcome ledger}: final selected institution ID and classification, relevant local instrument, official direct policy locator when available, search surfaces, negative-search closure date, and rationale for $\stP$ or $\stU$. A script may then reconstruct the six $c_j$ vectors, $D$, the institutional margins, and each sensitivity without manual edits.

Publication in a repository must preserve the difference between a file that was actually retrieved and a link that is currently discoverable. The provenance manifest should separately identify the public URL, the archived local digest, and whether redistribution of the content is permitted. When only metadata/link sharing is authorized, the source can be cited and hashed without republishing the publisher's PDF. Hashes should bind the record and exact edition; a successful recomputation of the arithmetic is not independent verification of the underlying classification. The accompanying CSV links all 57 institutional labels to the locators \emph{already embedded in the submitted six-page paper}; the existence of those links alone does not close the 36-record source-extraction or second-coder requirement.

An independent auditor should freeze and time-stamp all source objects before seeing the primary table, then classify them in randomized order, blind to the target jurisdiction's headline result. For any conflict, the adjudication log should cite the controlling clause and explain whether the disagreement concerns authority, scope, document status, timing, or language. Raw pre-adjudication counts are indispensable: an agreed-upon final label may conceal poor initial reproducibility. Agreement should be reported on the actual number of independently coded objects, not on duplicated translations or repeated records, and separately for layer admission and institutional-panel coding when their units differ. A revised DOI-bearing public release would then make the protocol citable and independently inspectable, rather than using the word ``reproducible'' as a substitute for a released ledger.

\section{Validity and Claim Boundaries}
The six-state corpus is finite, and individual jurisdictions are not sampled independent draws. Neither leave-one-jurisdiction-out checks nor the conditional perturbation bound are sampling confidence intervals. National direction, ministry HE instruments, and a single university rule have different institutional scales. The current audit tests \emph{availability of evidence at each scale}, not an actual directed chain of authority. A public gap can reflect source invisibility, archived web migration, language differences, or a legitimately decentralized control design. $\stU$ cannot be converted into ``no policy'' or ``noncompliant'' without different evidence.

The deterministic panel is bounded by public eligibility and founding-year information. Although it corrects two boundary omissions, the complete pre-selection universe was not preserved, limiting expanded oldest-$k$ sensitivity. Coder decisions were not independently double-coded; no inter-rater reliability, Arabic--English adjudication rate, search recall, or source-missingness distribution can be estimated from the supplied corpus. Interval-valued notation bounds only the permitted assignments to $\stP$ within the observed records. The published source dates, source identity, reproducible direct links and a non-identifying reviewer repository must therefore be checked before claiming independent reproducibility. The audit does not measure student behavior, enforcement, effective compliance, institutional quality, or policy impact.

\section{Conclusion}
At the frozen 23 July 2026 cutoff, the GCC exhibits six publicly explicit institution-level L4 existence cases but only one primary-gate explicit HE-specific L3 instrument. A broader UAE HE instrument gate raises L3 to two, while four jurisdictions still display the L3--L4 documentary discontinuity. The corrected oldest-four-with-ties panel identifies 20 explicit, 12 partial and 25 unlocated institutional cases; every jurisdiction contributes multiple explicit examples without supporting general prevalence. The robust contribution is a bounded and reusable \emph{public-evidence traceability protocol} that distinguishes instrument availability from documentary linkage and real implementation. It preserves the fuzzy envelope as transparent notation for partial observations and leaves independent source/coder validation as an explicit, necessary next step.

\begin{thebibliography}{25}\footnotesize
\bibitem{albous} M. R. Albous, O. R. Al-Jayyousi, and M. Stephens, ``AI governance in the GCC states: A comparative analysis of national AI strategies,'' \emph{Journal of Artificial Intelligence Research}, vol. 82, pp. 2389--2422, 2025.
\bibitem{jin} Y. Jin, L. Yan, V. Echeverria, D. Gasevic, and R. Martinez-Maldonado, ``Generative AI in higher education: A global perspective of institutional adoption policies and guidelines,'' \emph{Computers and Education: Artificial Intelligence}, vol. 8, p. 100348, 2025.
\bibitem{mcdonald} N. McDonald \emph{et al.}, ``Generative artificial intelligence in higher education: Evidence from an analysis of institutional policies and guidelines,'' \emph{Computers in Human Behavior: Artificial Humans}, vol. 3, p. 100121, 2025.
\bibitem{parker} L. Parker \emph{et al.}, ``Comparative analysis of artificial intelligence policies in universities across five countries,'' \emph{Discover Computing}, vol. 28, p. 267, 2025.
\bibitem{moorhouse} B. L. Moorhouse, M. A. Yeo, and Y. Wan, ``Generative AI tools and assessment: Guidelines of the world's top-ranking universities,'' \emph{Computers and Education Open}, vol. 5, p. 100151, 2023.
\bibitem{li} M. Li \emph{et al.}, ``A framework for developing university policies on generative AI governance: A cross-national comparative study,'' \emph{Studies in Higher Education}, 2026.
\bibitem{jiang} X. Jiang and Z. Abdullah, ``AI-enabled governance in higher education: A systematic review of applications, outcomes, and emerging implications,'' \emph{Frontiers in Education}, vol. 11, p. 1856440, 2026.
\bibitem{khoza} M. S. Khoza, S. F. Wamba, and S. Nyawa, ``Mapping the landscape of AI governance in higher education,'' \emph{Assessment \& Evaluation in Higher Education}, 2026.
\bibitem{an} Y. An, J. H. Yu, and S. James, ``Investigating higher education institutions' guidelines and policies regarding the use of generative AI in teaching, learning, research, and administration,'' \emph{International Journal of Educational Technology in Higher Education}, vol. 22, p. 10, 2025.
\bibitem{bacalso} F. D. Bacalso \emph{et al.}, ``Institutional approaches to generative AI management in higher education: A systematic review,'' \emph{Frontiers in Education}, vol. 11, p. 1814426, 2026.
\bibitem{qian} Y. Qian, ``Governing generative AI in higher education: Emerging policy approaches and support ecosystems at innovative U.S. universities,'' \emph{International Journal for Educational Integrity}, vol. 22, p. 25, 2026.
\bibitem{vukasovic} M. Vukasovic, ``Higher education governance and policy: An introduction to multi-issue, multi-level and multi-actor dynamics,'' \emph{Policy and Society}, vol. 36, no. 1, pp. 1--15, 2017.
\bibitem{donina} D. Donina and S. Hasanefendic, ``Higher education institutional governance reforms in the Netherlands, Portugal, and Italy: A policy translation perspective addressing the homogeneous/heterogeneous dilemma,'' \emph{Higher Education Quarterly}, vol. 73, no. 1, pp. 29--44, 2019.
\bibitem{squ} Sultan Qaboos University, ``Guidelines for using GenAI tools in teaching, learning and assessment,'' 2024. [Online]. Available: \url{https://www.squ.edu.om/About/Support-Centers-/Centre-For-Excellence-In-Teaching-and-Learning}.
\bibitem{obef} UAE Ministry of Higher Education and Scientific Research, ``University guidebook for outcome-based evaluation framework (OBEF), version 11.5,'' 23 Mar. 2026. [Online]. Available: \url{https://www.mohesr.gov.ae/Documents/OBF%20University%20Guide%20-%20EN.pdf}.
\bibitem{bqa1} Education and Training Quality Authority (BQA), Bahrain, ``Institutional review report: Kingdom University,'' 2019. [Online]. Available: \url{https://www.bqa.gov.bh/En/Reports/UniReports/HigherEducationReport/KU%20Inst%20en%20C2.pdf}.
\bibitem{bqa2} Education and Training Quality Authority (BQA), Bahrain, ``Higher education reports,'' 2026; index modified 7 July 2026. [Online]. Available: \url{https://www.bqa.gov.bh/institutions-performance/higher-education-reports/}.
\bibitem{caa1} Commission for Academic Accreditation, UAE, ``Renewal of licensure: Emirates Institute for Banking and Financial Studies---executive summary,'' 2019. [Online]. Available: \url{https://caa.ae/PORTALSUMMARYREPORTS/16/EIBFS%20Executive%20Summary.pdf}.
\bibitem{caa2} Commission for Academic Accreditation, UAE, ``Higher education institution: Emirates Institute for Banking and Financial Studies,'' 2026. [Online]. Available: \url{https://www.caa.ae/Pages/Institutes/Details.aspx?GUID=16}.
\bibitem{petry} F. E. Petry and R. R. Yager, ``Interval-valued fuzzy sets aggregation and evaluation approaches,'' \emph{Applied Soft Computing}, vol. 124, p. 108887, 2022.
\bibitem{pekala} B. P\k{e}kala, K. Dyczkowski, P. Grzegorzewski, and U. Bentkowska, ``Inclusion and similarity measures for interval-valued fuzzy sets based on aggregation and uncertainty assessment,'' \emph{Information Sciences}, vol. 547, pp. 1182--1200, 2021.
\bibitem{bahrain} Ministry of Education, Kingdom of Bahrain, ``Digital empowerment in education program,'' updated 30 Dec. 2025; school-sector target groups. [Online]. Available: \url{https://moe.gov.bh/en/digital-empowerment-in-education-program}.
\bibitem{unesco} UNESCO, \emph{Guidance for Generative AI in Education and Research}. Paris: UNESCO, 2023.
\bibitem{oecd} OECD, ``Policies supporting responsible and systematic GenAI adoption in higher education,'' \emph{OECD Education Spotlights}, no. 23, 2026.
\bibitem{qataru} Qatar University, ``Instructions on the use of generative artificial intelligence tools at Qatar University,'' approved/effective 25 Dec. 2025; English translation of Arabic original. [Online]. Available: \url{https://www.qu.edu.qa/en-us/about/policy/Documents/Instructions%20onThe%20Use%20of%20Generative%20Artificial%20Intelligence%20Tools%20at%20Qatar%20University-Approved.pdf}.
\end{thebibliography}
\end{document}
